\documentclass[11pt,a4paper]{article}
\usepackage[margin=1in]{geometry}
\usepackage{amsmath,amssymb,amsthm}
\usepackage{algorithm}
\usepackage{algpseudocode}
\usepackage{booktabs}
\usepackage{hyperref}

\theoremstyle{definition}
\newtheorem{definition}{Definition}[section]
\newtheorem{theorem}{Theorem}[section]
\newtheorem{lemma}[theorem]{Lemma}
\newtheorem{remark}{Remark}[section]

\title{\textbf{Theoretical Foundations, Expectation Preservation, and Architecture Ensembling of Stochastic Depth (DropPath)}}
\author{Team Scaibu \\ \small Email: \texttt{jaydeep.9664920749@gmail.com} \\ \small Architecture and Core Engine Team}
\date{October 2026}

\begin{document}
\maketitle

\begin{abstract}
This document provides the rigorous mathematical specification, expectation-invariance proofs, and effective-depth dynamics of Stochastic Depth (DropPath). In extremely deep residual networks and Vision Transformers, training time and gradient propagation decay with network depth. Stochastic depth randomly drops entire residual sub-blocks during training via Bernoulli gating while scaling surviving branches by $\frac{1}{1-p}$. We formulate the layer-level stochastic bypass, prove unbiased output expectation, derive computational complexity, trace a complete worked example, and analyze survival schedule design.
\end{abstract}

\section{Notation and Mathematical Preliminaries}

Let $x_l \in \mathbb{R}^{B \times D}$ denote the input activation tensor entering residual block $l \in \{1, \dots, L\}$, and let $F_l(x_l) \in \mathbb{R}^{B \times D}$ denote the transformation computed by the block's residual branch.

\begin{itemize}
    \item $B \in \mathbb{N}_{\ge 1}$: Batch size.
    \item $D \in \mathbb{N}_{\ge 1}$: Feature representation dimension.
    \item $L \in \mathbb{N}_{\ge 1}$: Total depth of residual blocks in the network.
    \item $p_l \in [0, 1)$: Drop probability for block $l$.
    \item $q_l = 1 - p_l \in (0, 1]$: Survival probability of block $l$.
    \item $b_{l, i} \sim \text{Bernoulli}(q_l)$: Binary survival gate variable for example $i \in \{1, \dots, B\}$ at block $l$.
    \item $x_{l+1} \in \mathbb{R}^{B \times D}$: Output tensor exiting the block.
\end{itemize}

\begin{table}[h]
\centering
\caption{Summary of Mathematical Symbols for Stochastic Depth}
\begin{tabular}{lll}
\toprule
\textbf{Symbol} & \textbf{Type / Domain} & \textbf{Description} \\
\midrule
$x_l$ & $\mathbb{R}^{B \times D}$ & Input residual tensor at layer $l$ \\
$F_l(x_l)$ & $\mathbb{R}^{B \times D}$ & Sublayer branch computation \\
$p_l$ & $[0, 1)$ & Layer drop probability \\
$q_l$ & $(0, 1]$ & Layer survival probability ($1 - p_l$) \\
$b_{l, i}$ & $\{0, 1\}$ & Per-sample Bernoulli survival indicator \\
$x_{l+1}$ & $\mathbb{R}^{B \times D}$ & Combined output tensor \\
\bottomrule
\end{tabular}
\end{table}

\section{Problem Definition and Core Concepts}

\subsection{Intuitive Meaning}
While standard Dropout zeroes individual neurons within a layer, Stochastic Depth drops entire transformation blocks ($F(x)$), allowing gradients to bypass the entire layer through the identity shortcut ($x + 0$). This dynamically samples networks of variable effective depth $\tilde{L} < L$ during training, alleviating vanishing gradients, shortening backward pass paths, and acting as an ensemble across exponentially many sub-architectures.

\subsection{Formal Mathematical Definition}
During training ($\text{training} = \text{true}$), for each sample $i \in \{1, \dots, B\}$:
\begin{align}
b_{l, i} &\sim \text{Bernoulli}(1 - p_l) \\
x_{l+1, i} &= x_{l, i} + \frac{b_{l, i}}{1 - p_l} F_l(x_{l, i})
\end{align}
During inference ($\text{training} = \text{false}$), the network evaluates all layers deterministically:
\begin{equation}
x_{l+1, i} = x_{l, i} + F_l(x_{l, i})
\end{equation}

\section{Algorithm Specification}

\begin{algorithm}[H]
\caption{Stochastic Depth (DropPath) Forward Execution}
\label{alg:stochastic_depth}
\begin{algorithmic}[1]
\Require Shortcut tensor $X \in \mathbb{R}^{B \times D}$, branch tensor $F \in \mathbb{R}^{B \times D}$, drop rate $p \in [0, 1)$, mode $\text{training}$, optional gates $b_{\text{in}} \in \{0, 1\}^B$.
\Ensure Combined output $Y \in \mathbb{R}^{B \times D}$, gate vector $b \in \{0, 1\}^B$, scale $s \in \mathbb{R}_{> 0}$.
\State $B \gets \text{rows}(X)$, $D \gets \text{cols}(X)$
\State Allocate $Y \in \mathbb{R}^{B \times D}$, $b \in \mathbb{R}^B$
\If{$\text{training} = \text{false} \lor p = 0$}
    \State $Y \gets X + F$, $b \gets \mathbf{1}^B$, $s \gets 1.0$
    \State \Return $\{Y, b, s\}$
\EndIf
\State $q \gets 1.0 - p$, $s \gets 1.0 / q$
\For{$i \gets 1 \textbf{ to } B$}
    \If{$b_{\text{in}}$ is provided}
        \State $b[i] \gets b_{\text{in}}[i]$
    \Else
        \State $u \sim \mathcal{U}(0, 1)$
        \State $b[i] \gets \mathbf{if } u < q \mathbf{ then } 1.0 \mathbf{ else } 0.0$
    \EndIf
    \For{$j \gets 1 \textbf{ to } D$}
        \State $Y_{ij} \gets X_{ij} + b[i] \cdot s \cdot F_{ij}$
    \EndFor
\EndFor
\State \Return $\{Y, b, s\}$
\end{algorithmic}
\end{algorithm}

\section{Theoretical Guarantees and Expectation Preservation}

\begin{theorem}[Unbiased Forward Expectation and Expected Depth]
For any fixed input $x_l$ and branch output $F_l(x_l)$:
\begin{equation}
\mathbb{E}_{b_l} [x_{l+1}] = x_l + F_l(x_l)
\end{equation}
and under a linear decay schedule $p_l = \frac{l}{L} p_{\text{max}}$, the expected effective depth $\mathbb{E}[\tilde{L}]$ of the network during training is:
\begin{equation}
\mathbb{E}[\tilde{L}] = L - \frac{p_{\text{max}} (L + 1)}{2}
\end{equation}
\end{theorem}

\begin{proof}
Taking the expectation with respect to the Bernoulli random variable $b_{l, i} \sim \text{Bernoulli}(1 - p_l)$:
\begin{equation}
\mathbb{E}[x_{l+1, i}] = \mathbb{E}\left[ x_{l, i} + \frac{b_{l, i}}{1 - p_l} F_l(x_{l, i}) \right] = x_{l, i} + \frac{\mathbb{E}[b_{l, i}]}{1 - p_l} F_l(x_{l, i}) = x_{l, i} + \frac{1 - p_l}{1 - p_l} F_l(x_{l, i}) = x_{l, i} + F_l(x_{l, i})
\end{equation}
The expected total number of active blocks is $\mathbb{E}[\tilde{L}] = \sum_{l=1}^L \mathbb{E}[b_l] = \sum_{l=1}^L (1 - p_l) = L - p_{\text{max}} \sum_{l=1}^L \frac{l}{L} = L - \frac{p_{\text{max}} L(L+1)}{2L} = L - \frac{p_{\text{max}}(L+1)}{2}$.
\end{proof}

\subsection{Limitations}
\begin{itemize}
    \item \textbf{High Early Drop Rates}: If early layers ($l \le 2$) have high drop rates ($p_l > 0.3$), feature extraction destabilizes before semantic representations form.
\end{itemize}

\section{Complexity Analysis}

\subsection{Time Complexity}
\begin{itemize}
    \item \textbf{Residual Addition}: Applying gating and addition across $B \times D$ elements requires $\mathcal{O}(B \cdot D)$ FLOPs.
    \item \textbf{Total Time Complexity}: $\Theta(B \cdot D)$ FLOPs.
\end{itemize}

\subsection{Space Complexity}
\begin{itemize}
    \item \textbf{Storage}: Requires $\mathcal{O}(B)$ floats to store per-sample gating indicators for backpropagation.
\end{itemize}

\section{Exactness and Error Analysis}

The scaling factor $\frac{1}{1-p}$ ensures that no post-hoc rescaling of model weights is required at test time, producing an algebraically exact expected transformation across all samples.

\section{Worked Numerical Example}

Let $B = 2, D = 2, p = 0.5 \implies q = 0.5, s = 2.0$.
\begin{equation}
X = \begin{bmatrix} 1.0 & 2.0 \\ 3.0 & 4.0 \end{bmatrix}, \quad F = \begin{bmatrix} 0.5 & 0.5 \\ 1.0 & 1.0 \end{bmatrix}, \quad b = [1.0, 0.0]^T
\end{equation}

\begin{enumerate}
    \item \textbf{Sample 1 (Survived, $b_1 = 1.0$)}:
    \begin{align}
    Y_{11} &= 1.0 + 1.0 \times 2.0 \times 0.5 = 1.0 + 1.0 = 2.0 \\
    Y_{12} &= 2.0 + 1.0 \times 2.0 \times 0.5 = 2.0 + 1.0 = 3.0
    \end{align}
    \item \textbf{Sample 2 (Dropped, $b_2 = 0.0$)}:
    \begin{align}
    Y_{21} &= 3.0 + 0.0 \times 2.0 \times 1.0 = 3.0 \\
    Y_{22} &= 4.0 + 0.0 \times 2.0 \times 1.0 = 4.0
    \end{align}
\end{enumerate}

\section{Numerical Considerations and Edge Cases}

\begin{itemize}
    \item \textbf{Drop Rate $p = 0$}: Identity residual connection $Y = X + F$.
    \item \textbf{Per-Sample vs Per-Batch Gating}: Sampling gates independently per token row ($b_i$) prevents entire mini-batches from skipping layers simultaneously, maintaining gradient flow consistency.
\end{itemize}

\begin{thebibliography}{9}
\bibitem{huang2016deep} G.~Huang, Y.~Sun, Z.~Liu, D.~Sedra, and K.~Q.~Weinberger, ``Deep Networks with Stochastic Depth,'' in \emph{European Conference on Computer Vision (ECCV)}, pp.~646--661, 2016.
\bibitem{touvron2021training} H.~Touvron et al., ``Training data-efficient image transformers & distillation through attention,'' in \emph{International Conference on Machine Learning (ICML)}, pp.~10347--10357, 2021.
\end{thebibliography}

\end{document}
